\textbf{Lombard \& Piraux 2004 --- Eq 38--42 (master ESIM equation, paper p.\ 15)}

\medskip

\textbf{Eq 38 --- assemble Eq 35 and Eq 37 into matrix
$\boldsymbol{M}$.}
Stacking one row per $(i, j) \in B = B_1 \cup B_2$ (rows on $B_1$
use Eq 35; rows on $B_2$ use Eq 37):
\begin{equation*}
  (\boldsymbol{U}^n)_B
  = \boldsymbol{M}
    \begin{pmatrix} \boldsymbol{W}_1^k \\ \boldsymbol{\Lambda}^k \end{pmatrix}
  + \begin{pmatrix} O(\Delta x^{k+1}) \\ \vdots \\ O(\Delta x^{k+1}) \end{pmatrix}
  \tag{38}
\end{equation*}
where $(\boldsymbol{U}^n)_B$ is the stacked vector of exact values
$\boldsymbol{U}(x_i, y_j, t^n)$ at the points of $B$. The radius
$q$ is chosen so that (38) is \emph{over-determined}.

\medskip

\textbf{Eq 39 --- least-squares pseudoinverse} $\boldsymbol{M}^{-1}$
(via normal equations, SVD, ...). Dropping the Taylor remainders
and identifying numerical estimates with exact values:
\begin{equation*}
  \begin{pmatrix} \boldsymbol{W}_1^k \\ \boldsymbol{\Lambda}^k \end{pmatrix}
  = \boldsymbol{M}^{-1} (\boldsymbol{U}^n)_B
  \tag{39}
\end{equation*}

\textbf{Eq 40 --- restrict to $\boldsymbol{W}_1^k$ only.} Let
$\bar{\boldsymbol{M}}^{-1}$ denote the restriction of
$\boldsymbol{M}^{-1}$ to the rows that produce
$\boldsymbol{W}_1^k$ (dropping the free-parameter rows for
$\boldsymbol{\Lambda}^k$):
\begin{equation*}
  \boldsymbol{W}_1^k = \bar{\boldsymbol{M}}^{-1} (\boldsymbol{U}^n)_B
  \tag{40}
\end{equation*}

\medskip

\textbf{Eq 41 --- modified value at $(I, J) \in \Omega_2$.} From
Eq 25, 26, 33:
\begin{equation*}
  (I, J) \in \Omega_2,\
  \boldsymbol{U}^*_{I,J}
  = \boldsymbol{\Pi}_{I,J}^k \boldsymbol{U}_1^k
  = \boldsymbol{\Pi}_{I,J}^k \boldsymbol{G}_1^k \boldsymbol{K}_1^k \boldsymbol{W}_1^k
  \tag{41}
\end{equation*}

\medskip

\textbf{Eq 42 --- master ESIM equation} (insert Eq 40 into Eq 41):
\begin{equation*}
  \boxed{\ \boldsymbol{U}^*_{I,J}
  = \boldsymbol{\Pi}_{I,J}^k \boldsymbol{G}_1^k \boldsymbol{K}_1^k
    \bar{\boldsymbol{M}}^{-1} (\boldsymbol{U}^n)_B\ }
  \tag{42}
\end{equation*}

\medskip

\textbf{Notes from the paper:}
\begin{itemize}
  \item For an irregular point on the $\Omega_1$ side,
    $\boldsymbol{G}_1^k$ and $\boldsymbol{K}_1^k$ are swapped for
    $\boldsymbol{G}_2^k$ and $\boldsymbol{K}_2^k$ in (42).
  \item For a fluid--solid interface ($\Omega_1$ = fluid,
    $\Omega_2$ = solid), the right-hand side of (42) involves
    BOTH the 3-component fluid solution AND the 5-component
    solid solution at the points of $B$.
\end{itemize}

\medskip

\textbf{Implementation notes (dipping-line Petrobras case):}
\begin{itemize}
  \item $\boldsymbol{\Pi}_{I,J}^k$, $\boldsymbol{G}_i^k$,
    $\boldsymbol{K}_i^k$, $\boldsymbol{S}_i^k$,
    $\bar{\boldsymbol{M}}^{-1}$ depend only on the relative
    positions $(x_i - x_P, y_j - y_P)$ and on the (constant)
    normal/tangent directions. They are pre-computed ONCE per
    irregular point during set-up.
  \item Per-timestep cost: one matrix--vector product
    $\bar{\boldsymbol{M}}^{-1} (\boldsymbol{U}^n)_B$ per
    irregular point, then a second product against
    $\boldsymbol{G}_1^k \boldsymbol{K}_1^k$, then a dot product
    with $\boldsymbol{\Pi}_{I,J}^k$. Total $\sim 1\%$ overhead at
    400$\times$400 per the paper's §3.9.
  \item For $k = 2$ on a fluid--solid interface, the paper's
    canonical radius is $q \approx 3.5\,\Delta x$ (gives
    $\sim 40$ neighbours so (38) is comfortably
    over-determined).
\end{itemize}
